% from ega2-1-fr.tex
\oldpage[II]{126}

On notera que cette dernière condition est automatiquement vérifiée lorsque
le préschéma localement noethérien $X$ est \emph{normal}
(\hyperref[II.6.4.9-fr]{6.4.9}).

D'autre part, l'hypothèse que
$\mathcal{B}\otimes_{\mathcal{A}}\mathcal{R}(X)$ soit localement libre peut
encore s'exprimer de la façon suivante~: désignons par $X_\alpha$ les
sous-préschémas fermés réduits de $X$ ayant pour espaces sous-jacents les
composantes irréductibles de $X$ (\hyperref[I.5.2.1-fr]{I, 5.2.1}), qui sont
donc des préschémas intègres localement noethériens. Tout $x\in X$ appartient
à un nombre fini des sous-espaces $X_\alpha$~; d'autre part,
$\mathcal{B}\otimes_{\mathcal{A}}\mathcal{R}(X_\alpha)$ est un
$\mathcal{R}(X_\alpha)$-Module localement libre de rang constant $k_\alpha$
(\hyperref[I.7.3.6-fr]{I, 7.3.6})~; dire que
$\mathcal{B}\otimes_{\mathcal{A}}\mathcal{R}(X)$ est un
$\mathcal{R}(X)$-Module localement libre signifie que, pour tout $x\in X$,
\emph{les rangs $k_\alpha$ correspondant aux indices tels que $x\in X_\alpha$
sont égaux}. La question est en effet locale, et on est ramené au cas
$X=\operatorname{Spec}(C)$, où $C$ est un anneau noethérien réduit, et
$\mathcal{B}=\widetilde{D}$, où $D$ est une $C$-algèbre qui est un $C$-module
de type fini~; si $\mathfrak{p}_i$ ($1\leq i\leq m$) sont les idéaux premiers
minimaux de $C$, l'anneau total des fractions $L$ de $C$ est composé direct
des corps des fractions $K_i$ des anneaux intègres $C/\mathfrak{p}_i$, et
$D\otimes_C L$ est somme directe des $D\otimes_C K_i$, d'où la conclusion.

Il est clair que sous les hypothèses (I) ou (II), on a défini ainsi un
\emph{homomorphisme de faisceaux de monoïdes multiplicatifs}
$N_{\mathcal{B}/\mathcal{A}}:\mathcal{B}\to\mathcal{A}$, aussi noté $N$ si
aucune confusion n'en résulte, et appelé l'homomorphisme \emph{norme}. Pour
deux sections $f$, $g$ de $\mathcal{B}$ au-dessus d'un même ouvert $U$, on a
donc
\[
\label{II.6.5.1.1-fr}
N_{\mathcal{B}/\mathcal{A}}(fg)
=N_{\mathcal{B}/\mathcal{A}}(f)N_{\mathcal{B}/\mathcal{A}}(g)
\tag{6.5.1.1}
\]
pour les sections correspondantes de $\mathcal{A}$ au-dessus de $U$~;
\[
\label{II.6.5.1.2-fr}
N_{\mathcal{B}/\mathcal{A}}(1_{\mathcal{B}})=1_{\mathcal{A}}~;
\tag{6.5.1.2}
\]
enfin, pour toute section $s$ de $\mathcal{A}$ au-dessus de $U$
\[
\label{II.6.5.1.3-fr}
N_{\mathcal{B}/\mathcal{A}}(s\cdot1_{\mathcal{B}})=s^n
\tag{6.5.1.3}
\]
si le rang de $\mathcal{B}$ est constant et égal à $n$ (pour
$s=0_{\mathcal{A}}$, cette formule donne
$N(0_{\mathcal{B}})=0_{\mathcal{A}}$ si $n\geq1$,
$N(0_{\mathcal{B}})=N(1_{\mathcal{B}})=1_{\mathcal{A}}$ si $n=0$).

Dans l'hypothèse (I), pour que $f\in\Gamma(U,\mathcal{B})$ soit inversible,
il faut et il suffit que $N(f)\in\Gamma(U,\mathcal{A})$ le soit. Dans
l'hypothèse (II), cette condition est nécessaire~; elle est suffisante (par
(\hyperref[II.6.4.7-fr]{6.4.7})) lorsqu'on suppose que
$\mathcal{B}\to\mathcal{B}\otimes_{\mathcal{A}}\mathcal{R}(X)$ est
\emph{injectif} et que l'hypothèse plus restrictive suivante est vérifiée~:
\begin{enumerate}
\item[(II \emph{bis})] \emph{$(X,\mathcal{A})$ est un préschéma localement
noethérien réduit, $\mathcal{B}$ est un $\mathcal{A}$-Module cohérent tel que
$\mathcal{B}\otimes_{\mathcal{A}}\mathcal{R}(X)$ soit un
$\mathcal{R}(X)$-Module localement libre, et pour toute section
$f\in\Gamma(U,\mathcal{B})$ au-dessus d'un ouvert tel que
$\mathcal{B}\otimes_{\mathcal{A}}\mathcal{R}(X)|U$ soit de rang constant $n$
sur $\mathcal{R}(X)|U$, les $\sigma_j(f)$ ($1\leq j\leq n$) sont des sections
de $\mathcal{A}$ au-dessus de $U$.}
\end{enumerate}
(On notera encore que cette condition est vérifiée lorsque $X$ est
\emph{normal}.)
\end{env}

\begin{env}[6.5.2]
\label{II.6.5.2-fr}
Supposons vérifiée une des hypothèses (I), (II) de
(\hyperref[II.6.5.1-fr]{6.5.1}), et soit $\mathcal{L}'$ un
$\mathcal{B}$-Module \emph{inversible}. Nous allons lui associer canoniquement
(à un isomorphisme unique près) un $\mathcal{A}$-Module \emph{inversible}
$\mathcal{L}$ de la façon suivante. Désignons par $\mathcal{A}^*$ (resp.
$\mathcal{B}^*$) le sous-faisceau de $\mathcal{A}$ (resp. $\mathcal{B}$) tel
que $\Gamma(U,\mathcal{A}^*)$ (resp. $\Gamma(U,\mathcal{B}^*)$) soit
l'ensemble des

